% !PRINTED 377
\weberpage{377}
Ist umgekehrt die Kongruenz (3) erfüllt, so kann man
$D = n^2 - 4\,ml$ setzen und erhält eine Form der Diskriminante $D$,
$(m, n, l)$, durch die $m$ eigentlich darstellbar ist.

1. \gesperrt{Die notwendige und hinreichende Bedingung dafür, daß eine zu $Q$ teilerfremde Zahl $m$ durch eine Form der Diskriminante $D$ eigentlich darstellbar ist, besteht also darin, daß die Kongruenz}

\weq{4}{x^2 \equiv D \pmod{4m}}

\gesperrt{lösbar sei.}

Aus einer Form $(a,b,c)$ erhält man eine Schar \gesperrt{äquivalenter
Formen}:

\weq{5}{(a,\ b+2\lambda a,\ C),}

worin $\lambda$ eine beliebige ganze Zahl und $C$ durch die Gleichung

\[ D = (b+2\lambda a)^2 - 4aC \]

bestimmt ist und sich gleich $a\lambda^2+b\lambda+c$ ergibt. Das System (5)
wird eine \gesperrt{Schar paralleler Formen} genannt.

Wenn die Kongruenz (4) erfüllt ist, so ist auch

\[ (x+2\lambda m)^2 \equiv D \pmod{4m}, \]

und wir wollen die ganze Schar der Wurzeln von (4), die nach
dem Modul $2m$ kongruent sind, als \gesperrt{eine Wurzel} betrachten. In
diesem Sinne sei die Anzahl der Wurzeln von (4), die demnach
sicher endlich ist, mit $\psi(D,m)$ zu bezeichnen.

2. \gesperrt{Jede Wurzel von (4) gibt Anlaß zu einer Schar paralleler Formen, durch die die Zahl $m$ eigentlich darstellbar ist.}

Die Zahl $\psi(D,m)$ läßt sich leicht bestimmen. Sind zunächst
$m'$, $m''$ relativ prim, so ist

\weq{6}{\psi(D,m'),\ \psi(D,m'') = \psi(D,m'm'').} % [sic] Gedruckt ist ein Komma statt eines Multiplikationszeichens.

Denn ist

\weq{7}{x'^2 \equiv D \pmod{4m'}, \quad x''^2 \equiv D \pmod{4m''},}

und

\weq{8}{\begin{aligned} x &\equiv x' \pmod{2m'}, \\ &\equiv x'' \pmod{2m''}, \end{aligned}}

so ist $(x^2-D):4$ durch $m'$ und durch $m''$, also auch durch $m'm''$
teilbar, und die Kongruenz

\weq{9}{x^2 \equiv D \pmod{4m'm''}}

erfüllt. Umgekehrt gibt jede Wurzel der Kongruenz (9) je eine
Wurzel der beiden Kongruenzen (7). Also erhält man alle
%CONT>>
